\documentstyle{article}
\begin{document}

\centerline{\bf Equivalence Classes and Symmetries of Abel's Equation}

\vskip 0.3cm

\centerline{Fritz Schwarz}
\centerline{GMD, Institut SCAI}
\centerline{53754 Sankt Augustin, Germany}
\centerline{Email: fritz.schwarz@gmd.de}

\vskip 0.3cm

A scheme for solving Abel's equation is described that is based on a
certain infinite group of point transformations, the so-called {\em
structure invariance} group, and a normal form wrt. this group, the
so-called {\em rational normal form (RNF)}. The key objects in this
scheme are the {\em equivalence classes} of equations wrt. this group.
They are identified by means of an {\em absolute invariant}. The
equations of an equivalence class may allow one or two point symmetries
that are applied for the solution procedure. If an equation does not
have a symmetry, membership in an equivalence class may be decided by
functional decomposition in the coefficient field of the given
equation. Based on this scheme a solution algorithm is developed that
generates most of the closed form solutions described in the literature
as a consequence of non-trivial symmetries. It is suggested that
similar schemes may be developed for other classes of equations.

\end{document}
